% Options for packages loaded elsewhere
\PassOptionsToPackage{unicode}{hyperref}
\PassOptionsToPackage{hyphens}{url}
\PassOptionsToPackage{dvipsnames,svgnames,x11names}{xcolor}
\documentclass[
  french,
  11pt,
  a4paper,
]{article}
\usepackage{xcolor}
\usepackage[margin=2.3cm]{geometry}
\usepackage{amsmath,amssymb}
\setcounter{secnumdepth}{5}
\usepackage{iftex}
\ifPDFTeX
  \usepackage[T1]{fontenc}
  \usepackage[utf8]{inputenc}
  \usepackage{textcomp} % provide euro and other symbols
\else % if luatex or xetex
  \usepackage{unicode-math} % this also loads fontspec
  \defaultfontfeatures{Scale=MatchLowercase}
  \defaultfontfeatures[\rmfamily]{Ligatures=TeX,Scale=1}
\fi
\usepackage{lmodern}
\ifPDFTeX\else
  % xetex/luatex font selection
  \setmainfont[]{TeX Gyre Pagella}
  \setmathfont[]{TeX Gyre Pagella Math}
\fi
% Use upquote if available, for straight quotes in verbatim environments
\IfFileExists{upquote.sty}{\usepackage{upquote}}{}
\IfFileExists{microtype.sty}{% use microtype if available
  \usepackage[]{microtype}
  \UseMicrotypeSet[protrusion]{basicmath} % disable protrusion for tt fonts
}{}
\makeatletter
\@ifundefined{KOMAClassName}{% if non-KOMA class
  \IfFileExists{parskip.sty}{%
    \usepackage{parskip}
  }{% else
    \setlength{\parindent}{0pt}
    \setlength{\parskip}{6pt plus 2pt minus 1pt}}
}{% if KOMA class
  \KOMAoptions{parskip=half}}
\makeatother
\usepackage{longtable,booktabs,array}
\newcounter{none} % for unnumbered tables
\usepackage{calc} % for calculating minipage widths
% Correct order of tables after \paragraph or \subparagraph
\usepackage{etoolbox}
\makeatletter
\patchcmd\longtable{\par}{\if@noskipsec\mbox{}\fi\par}{}{}
\makeatother
% Allow footnotes in longtable head/foot
\IfFileExists{footnotehyper.sty}{\usepackage{footnotehyper}}{\usepackage{footnote}}
\makesavenoteenv{longtable}
\ifLuaTeX
\usepackage[bidi=basic,shorthands=off]{babel}
\else
\usepackage[bidi=default,shorthands=off]{babel}
\fi
\ifPDFTeX
\else
\babelfont{rm}[]{TeX Gyre Pagella}
\fi
\ifLuaTeX
  \usepackage{selnolig} % disable illegal ligatures
\fi
\setlength{\emergencystretch}{3em} % prevent overfull lines
\providecommand{\tightlist}{%
  \setlength{\itemsep}{0pt}\setlength{\parskip}{0pt}}
\usepackage[style=apa]{biblatex}
\addbibresource{../../02_ETAT_DE_L_ART_FINAL/references.bib}
\usepackage{etoolbox,caption}\captionsetup{font=small,labelfont=bf,labelsep=endash}\captionsetup[table]{position=top}\AtBeginEnvironment{longtable}{\small}\addto\captionsfrench{\renewcommand{\tablename}{Tableau}}
\usepackage{bookmark}
\IfFileExists{xurl.sty}{\usepackage{xurl}}{} % add URL line breaks if available
\urlstyle{same}
\hypersetup{
  pdftitle={Perspectives de recherche --- RL-in-GA pour le problème du sac à dos et au-delà},
  pdfauthor={RALAIVAO Niaiko Michaël --- Master Recherche, ENI / LIMAD / GLoRIA},
  pdflang={fr},
  colorlinks=true,
  linkcolor={blue},
  filecolor={Maroon},
  citecolor={teal},
  urlcolor={blue},
  pdfcreator={LaTeX via pandoc}}

\title{Perspectives de recherche --- RL-in-GA pour le problème du sac à
dos et au-delà}
\author{RALAIVAO Niaiko Michaël --- Master Recherche, ENI / LIMAD /
GLoRIA}
\date{Septembre 2026}

\begin{document}
\maketitle

{
\setcounter{tocdepth}{3}
\tableofcontents
}
Ce document recense ce qui \textbf{manque} au travail actuel, les
\textbf{extensions} immédiates, les \textbf{interactions avec d'autres
domaines} et les pistes de \textbf{généralisation}. Les pistes sont
classées par horizon (court : quelques semaines ; moyen : un semestre ;
long : une thèse).

\section{Ce qui manque au travail actuel (limites à lever en
priorité)}\label{perspectives-ce-qui-manque-au-travail-actuel-limites-uxe0-lever-en-priorituxe9}

{\def\LTcaptype{none} % do not increment counter
\begin{longtable}[]{@{}
  >{\raggedright\arraybackslash}p{(\linewidth - 6\tabcolsep) * \real{0.2500}}
  >{\raggedright\arraybackslash}p{(\linewidth - 6\tabcolsep) * \real{0.2500}}
  >{\raggedright\arraybackslash}p{(\linewidth - 6\tabcolsep) * \real{0.2500}}
  >{\raggedright\arraybackslash}p{(\linewidth - 6\tabcolsep) * \real{0.2500}}@{}}
\toprule\noalign{}
\begin{minipage}[b]{\linewidth}\raggedright
Limite
\end{minipage} & \begin{minipage}[b]{\linewidth}\raggedright
Conséquence
\end{minipage} & \begin{minipage}[b]{\linewidth}\raggedright
Action proposée
\end{minipage} & \begin{minipage}[b]{\linewidth}\raggedright
Horizon
\end{minipage} \\
\midrule\noalign{}
\endhead
\bottomrule\noalign{}
\endlastfoot
Population initiale aléatoire et budget indépendant de n & Au-delà de 5
000 objets, le glouton bat l'AG (E5) & Initialisation gloutonne d'une
fraction de la population (option \texttt{init\_greedy\_fraction} déjà
codée) ; budget proportionnel à n & court \\
Budget d'évaluation unique (P = 100, G = 200) & Les conclusions peuvent
dépendre du budget : avec plus de générations, le QL en ligne aurait
plus de temps pour apprendre & Courbes « qualité vs budget » (G ∈ \{50,
100, 200, 500, 1 000\}) ; comparaison à temps de calcul égal & court \\
Une seule métaheuristique de base et une seule réparation (très
efficace) & La réparation gloutonne « lisse » les différences entre
opérateurs & Répéter avec pénalisation au lieu de réparation ; faire de
la stratégie de réparation une \textbf{action} (GAP 1, action \(r_t\)) &
court \\
Signal de récompense qui s'éteint en fin de recherche & Les états « fin
de recherche » sont peu ou pas appris (Q ≈ 0) & Récompenses normalisées
par rang ou relatives aux améliorations récentes ; attribution de crédit
par valeurs extrêmes (EVCA) ou fenêtre glissante
\autocite{karimi-mamaghan_machine_2022} & court \\
Espace d'état discret et grossier (27 états) & Généralisation limitée ;
seuils calibrés à la main & Q-learning avec approximation de fonction
(tuile, MLP) ; \textbf{DQN} \autocite{mnih_2015} ou \textbf{PPO}
\autocite{schulman_2017} sur l'état continu & moyen \\
Descripteur d'instance φ(I) scalaire & Les états « classe difficile » ne
sont vus qu'à l'entraînement & Représentation apprise de l'instance et
de la population : Deep Sets / \textbf{Set Transformer}
\autocite{lee_settransformer_2019} sur les objets \((w_i, v_i, p_i)\)
(GAP 3) & moyen \\
10 exécutions par instance, instances de taille ≤ 2 000 (sauf E5) &
Puissance statistique limitée sur les petites différences & 30
exécutions ; instances de Pisinger « hard » (déjà téléchargées) et 3 240
instances de Jooken & court \\
Comparaison exacte limitée à la PD et à HiGHS & Pas de comparaison avec
\emph{Combo} ou \emph{Expknap} & Intégrer \emph{Combo} (code C de
Pisinger) comme référence exacte spécialisée & court \\
\end{longtable}
}

\section{Extensions
méthodologiques}\label{perspectives-extensions-muxe9thodologiques}

\begin{enumerate}
\def\labelenumi{\arabic{enumi}.}
\tightlist
\item
  \textbf{Contrôle continu des paramètres} : au lieu de 9 actions
  discrètes, l'agent règle directement \((p_c, p_m, k)\) --- politique
  continue (PPO, SAC) ; comparaison avec la configuration dynamique
  d'algorithmes \autocite{biedenkapp_2020}.
\item
  \textbf{Méta-apprentissage et transfert systématique} : entraîner sur
  une distribution hétérogène d'instances et mesurer la généralisation
  hors distribution selon trois axes --- taille, corrélation, serrage
  --- comme proposé dans l'hypothèse GAP 1 ; sélection d'instances
  d'entraînement représentatives \autocite{benjamins_2024}.
\item
  \textbf{Portefeuille de politiques (GAP 2)} : une archive MAP-Elites
  \autocite{mouret_clune_2015} indexée par les caractéristiques
  d'instance (corrélation × serrage) contenant des politiques
  spécialisées, sélectionnées à l'exécution.
\item
  \textbf{Boucle fermée AR ↔ AG (type EAM)} : une politique neuronale
  constructive initialise la population et l'AG affine ; les solutions
  évoluées ré-entraînent la politique \autocite{gu_eam_2025} --- à
  étendre au KP et à l'évaluation hors distribution.
\item
  \textbf{Hyper-heuristique} : l'agent choisit entre heuristiques
  complètes (AG, recuit simulé, recherche tabou, recherche locale)
  plutôt qu'entre opérateurs
  \autocite{kallestad_general_2023,dokeroglu_hyper-heuristics_2024}.
\item
  \textbf{Cadre théorique} : relier la diversité de population (ou
  l'entropie des fréquences \(p_i\)) à la capacité de généralisation
  (question QT1/QT4 du brainstorming) ; borne sur le biais introduit par
  l'agent (à la manière de la borne KL d'EAM).
\end{enumerate}

\section{Autres variantes du sac à dos et autres
problèmes}\label{perspectives-autres-variantes-du-sac-uxe0-dos-et-autres-probluxe8mes}

{\def\LTcaptype{none} % do not increment counter
\begin{longtable}[]{@{}
  >{\raggedright\arraybackslash}p{(\linewidth - 4\tabcolsep) * \real{0.3333}}
  >{\raggedright\arraybackslash}p{(\linewidth - 4\tabcolsep) * \real{0.3333}}
  >{\raggedright\arraybackslash}p{(\linewidth - 4\tabcolsep) * \real{0.3333}}@{}}
\toprule\noalign{}
\begin{minipage}[b]{\linewidth}\raggedright
Cible
\end{minipage} & \begin{minipage}[b]{\linewidth}\raggedright
Ce qui change
\end{minipage} & \begin{minipage}[b]{\linewidth}\raggedright
Intérêt
\end{minipage} \\
\midrule\noalign{}
\endhead
\bottomrule\noalign{}
\endlastfoot
MKP de grande taille (GK, OR30x500) & rien dans l'architecture &
NP-difficile au sens fort : la PLNE atteint vite ses limites \\
Sac à dos avec graphe de conflits (KPCG)
\autocite{pferschy_knapsack_2009}, avec précédences (PCKP) & réparation
respectant les conflits / précédences & lien avec la théorie des graphes
et la coloration \\
Sac à dos quadratique, multiple, avec \emph{setups}
\autocite{cacchiani_knapsack_2022} & évaluation non linéaire, codage
entier & variantes sans algorithme dédié performant \\
KP stochastique, dynamique, \emph{chance-constrained} & l'état inclut
l'incertitude ; ré-optimisation & logistique en temps réel \\
KP multi-objectif (valeur vs risque) & NSGA-II \autocite{deb_fast_2002}
piloté par AR & aide à la décision multicritère (axe du LIMAD) \\
TSP / CVRP, ordonnancement (FJSP) & codage en permutation, opérateurs OX
/ 2-opt & valider la \textbf{généralité} du cadre (GAP 1, cas TSP)
\autocite{chen_slga_2020} \\
\end{longtable}
}

\section{Interactions avec d'autres
domaines}\label{perspectives-interactions-avec-dautres-domaines}

\begin{itemize}
\tightlist
\item
  \textbf{Aide à la décision territoriale et SIG} (LIMAD) : sélection de
  sites d'équipements, planification de projets communaux géolocalisés
  --- prolongement direct de la communication de mai 2024 (cartographie
  numérique).
\item
  \textbf{Logistique humanitaire} (gestion des risques et catastrophes)
  : chargement multi-contraintes, affectation de véhicules, avec données
  réelles d'organismes partenaires.
\item
  \textbf{Infonuagique et réseaux} : placement de machines virtuelles,
  allocation de bande passante --- problèmes MKP dynamiques.
\item
  \textbf{Finance} : portefeuilles sous contraintes (KP moyenne-VaR
  \autocite{vaezi_mean-var_2020}).
\item
  \textbf{Génie logiciel} (équipe GLoRIA) : sélection de fonctionnalités
  à livrer sous contrainte de budget (\emph{next release problem},
  variante du KP), sélection de tests de régression, allocation de
  ressources dans les projets agiles.
\item
  \textbf{Modèles de langage} : génération ou sélection d'opérateurs par
  un LLM, avec l'AR pour arbitrer --- piste émergente.
\end{itemize}

\section{Feuille de route proposée (vers une
thèse)}\label{perspectives-feuille-de-route-proposuxe9e-vers-une-thuxe8se}

{\def\LTcaptype{none} % do not increment counter
\begin{longtable}[]{@{}
  >{\raggedright\arraybackslash}p{(\linewidth - 4\tabcolsep) * \real{0.3333}}
  >{\raggedright\arraybackslash}p{(\linewidth - 4\tabcolsep) * \real{0.3333}}
  >{\raggedright\arraybackslash}p{(\linewidth - 4\tabcolsep) * \real{0.3333}}@{}}
\toprule\noalign{}
\begin{minipage}[b]{\linewidth}\raggedright
Phase
\end{minipage} & \begin{minipage}[b]{\linewidth}\raggedright
Contenu
\end{minipage} & \begin{minipage}[b]{\linewidth}\raggedright
Valorisation visée
\end{minipage} \\
\midrule\noalign{}
\endhead
\bottomrule\noalign{}
\endlastfoot
1 (3 mois) & Levée des limites courtes : budget, réparation comme
action, récompense normalisée, 30 exécutions, instances \emph{hard} &
article de conférence (CARI, GECCO \emph{student workshop}) \\
2 (6 mois) & État continu + DQN/PPO ; représentation Set Transformer ;
protocole OOD complet (taille × corrélation × serrage) & article
(EvoCOP, LION) \\
3 (6--12 mois) & Généralisation TSP/FJSP ; portefeuille de politiques ;
cas réels avec partenaires & revue (C\&OR, EJOR, \emph{Applied Soft
Computing}) \\
\end{longtable}
}

\printbibliography[title=Bibliographie]

\end{document}
